% 本文件由 03_代码/p4_tables.py 自动生成，请勿手工修改。
% 数据源：06_支撑材料/p4_results.json

\begin{table}[htbp]
  \centering
  \caption{两个交付文件在 2025 年 2 月 1 日--12 月 31 日（334 天）的实际运行结果}
  \label{tab:p4-branch}
  \small
  \begin{tabular}{lrr}
    \toprule
    指标 & result4-2（计划不得修改） & result4-3（6/12/18 可调整） \\
    \midrule
    计划量（$0{:}00$ 定下的量）/kWh & 21,705,435 & 21,399,492 \\
    最终生效常规量/kWh & 21,705,435 & 20,842,200 \\
    上调量/kWh & --- & 436,206 \\
    下调量/kWh & --- & 993,498 \\
    紧急购电量/kWh & 116,391 & 29,325 \\
    \midrule
    弃电量/kWh & 2,645,486 & 1,667,712 \\
    计划费/元 & 14,041,884 & 13,802,989 \\
    调整费/元 & --- & 197,228 \\
    紧急费/元 & 698,929 & 153,069 \\
    合计费用/元 & \textbf{14,740,813} & \textbf{14,153,286} \\
    期初储电量/kWh & 7,007 & 6,827 \\
    期末储电量/kWh & 6,242 & 5,890 \\
    \bottomrule
  \end{tabular}
  \par\vspace{2pt}
  \begin{minipage}{\textwidth}\footnotesize
    注：两个分支均在同一条附件 4 实际价格路径上分别进行时序仿真和参数标定，费用分别统计。两者均按对应时段的实际电价结算，并包含初始计划费。4-3 主口径允许上调和下调，调减部分退回原价并按 $50\%$ 的实际电价支付违约费；4-2 的全天计划固定，不发生调整。
  \end{minipage}
\end{table}
